\chapter{Every Command}\label{cha:commands}

This chapter is the whole vocabulary of the prompt, one line each, and it
is not written by hand: \pth{doc/guide/mkref.py} reads the shell's own
command table out of the ROM's source, the programs out of
\pth{/SYSTEM/BIN} and the languages out of \pth{/LANG} at the moment
this book is built. A command the machine has and this list does not
describe stops the build, and so does a description of a command that has
gone --- so what is below is what the machine you are reading about
answers to. Chapter~\ref{cha:shell} is how the shell behaves; this is what
you can say to it.

Each entry says where the word lives, which matters more than it looks:

\begin{description}[leftmargin=1.4em,style=nextline]
\item[ROM] in the resident part of the operating system, always there.
\item[ROM, bank \emph{n}] in one of the sideways banks
  (Chapter~\ref{cha:memory}), paged in for the moment it runs.
\item[runs MONITOR, runs TYPE] a word the ROM still answers to, which
  hands its line to a program on the disk. The work moved out of the ROM
  in September~2026 to make room; the word did not change.
\item[/SYSTEM/BIN] a program. Programs run by name from anywhere, and
  take their arguments the same way the ROM's words do.
\item[/LANG/\emph{name}] a language: its folder, with its examples.
\item[ROM, on the Linux] a word the ROM answers to by starting something
  on the Linux the machine runs on (Chapter~\ref{cha:linux}).
\end{description}

\noindent Every word below also works from a BASIC, Forth or the monitor
with a \verb|*| in front, and a bare name that is not here is tried as a
program, then as a CP/M program if F7 allows it, then as an alias, then as
an RX script --- in that order, so an alias never hides a command.
Capitals are the convention; the shell does not mind either way.

\input{generated/commands}
